% Originally: the \section{Solutions} of content/week-03.tex

\section{Solutions}
\label{sec:compactness_solutions}

\begin{exercisesolution}[Solution to \cref{ex:classification_table}]
\label{sol:classification_table}
% Generator: Claude Opus 5 (Medium)
\begin{center}
\begin{tabular}{@{}l cccc@{}}
\textbf{Set} & \textbf{open} & \textbf{closed} & \textbf{compact} & \textbf{complete} \\
\hline
$[0,1]$ & $\times$ & $\checkmark$ & $\checkmark$ & $\checkmark$ \\
$\mathbb{Q}$ & $\times$ & $\times$ & $\times$ & $\times$ \\
$B_1(0) \subseteq \mathbb{R}^3$ & $\checkmark$ & $\times$ & $\times$ & $\times$ \\
$\left\{\tfrac1n\right\}$ & $\times$ & $\times$ & $\times$ & $\times$ \\
$\{0\}\cup\left\{\tfrac1n\right\}$ & $\times$ & $\checkmark$ & $\checkmark$ & $\checkmark$ \\
$\bigcup_{n\geq1}B_{1/n}(n)$ & $\checkmark$ & $\times$ & $\times$ & $\times$ \\
$(-\infty,1]$ & $\times$ & $\checkmark$ & $\times$ & $\checkmark$ \\
$f([0,1])$ & $\times$ & $\checkmark$ & $\checkmark$ & $\checkmark$ \\
$f^{-1}([0,1])$ & \textbf{?} & $\checkmark$ & \textbf{?} & $\checkmark$ \\
$f^{-1}\big((0,1)\big)$ & $\checkmark$ & \textbf{?} & \textbf{?} & \textbf{?} \\
$S = \operatorname{span}\{(1,0),(3,0)\}$ & $\times$ & $\checkmark$ & $\times$ & $\checkmark$ \\
\end{tabular}
\end{center}

The rows that carry the lesson:

\textbf{$\left\{\tfrac1n\right\}$ versus $\{0\}\cup\left\{\tfrac1n\right\}$.} Adding the single
point $0$ changes all four answers. Without it the set misses its own limit point, so it is not
closed, and the Cauchy sequence $\tfrac1n$ has nowhere to converge inside the set. With it, the
set is closed and bounded, hence compact by \cref{thm:heine_borel}.

\textbf{$\bigcup_{n\geq1}B_{1/n}(n)$.} Open as a union of open sets
(\cref{item:arbitrary_union_open}); not closed, since each endpoint $n\pm\tfrac1n$ is a limit
point outside the set; unbounded, so certainly not compact.

\textbf{$S$ is a trap.} The two spanning vectors are parallel, so $S$ is not $\mathbb{R}^2$ but
the $x$-axis, a line, which is closed, unbounded and not compact.

\textbf{The $f$ rows show exactly how far continuity gets you.} Preimages behave: $f^{-1}$ of a
closed set is closed, of an open set is open (\cref{def:continuous}). Images do not, and only
compactness transfers forward (\cref{item:continuity_preserves_compact}), which is why
$f([0,1])$ is compact but $f^{-1}([0,1])$ need not be: take $f \equiv 0$, giving
$f^{-1}([0,1]) = \mathbb{R}^n$. That same $f$ makes $f^{-1}([0,1])$ open, while $f = \id$ makes
it $[0,1]$, which is not, hence the \textbf{?}. And $f([0,1])$ is never open: it is compact,
hence closed and bounded, and a non-empty proper subset of $\mathbb{R}^m$ cannot be both open
and closed, since $\mathbb{R}^m$ is connected (\cref{def:connectedness}).

% Reclassified from ainote to remark: an observation about the table's mathematical content.
\begin{remark}[Why the last two columns agree]
The punchline of the table is that its last two columns are \textbf{identical}. That is not a
coincidence of this particular list. Inside a \emph{complete} ambient space a subset is complete
exactly when it is closed (\cref{rem:complete_vs_closed}), and $\mathbb{R}^n$ is complete. Build
the same table inside $\mathbb{Q}$ instead and the two columns come apart immediately.
\end{remark}
\end{exercisesolution}

% Transition: Claude Opus 5 (High)
% Correction: Claude Opus 5 (High) --- replaces an ainote that was wrong three ways. It named
% ex:connectedness_tf as being "in this chapter" when that exercise is in ch. 8; it claimed only
% aiexercise problems are solved in this section, though ex:open_closed_complete_compact and ex:classification_table are
% here and are neither; and it listed two of the six exercises actually solved inline.
Not every solution in this chapter is collected here. \Cref{ex:totally_bounded_implies_bounded}
(\cpageref{ex:totally_bounded_implies_bounded}), \cref{ex:compact_finite_complete}
(\cpageref{ex:compact_finite_complete}), \cref{ex:convergent_sequence_compact}
(\cpageref{ex:convergent_sequence_compact}), \cref{ex:finite_set_compact}
(\cpageref{ex:finite_set_compact}), \cref{ex:topological_properties_table}
(\cpageref{ex:topological_properties_table}) and \cref{ex:heine_borel_fails}
(\cpageref{ex:heine_borel_fails}) are solved where they appear.

\begin{exercisesolution}[Solution to \cref{ex:open_closed_complete_compact}]
\label{sol:open_closed_complete_compact}
% Generator: Claude Opus 5 (High)
The efficient route is to write each set as a preimage of an open or closed set under a
continuous map, and then read openness or closedness off \cref{def:continuous}\textbf{(c)}
rather than chasing balls. Compactness is then Heine--Borel (\cref{thm:heine_borel}): closed
\emph{and} bounded.
\begin{center}
\begin{tabular}{@{}c l c c c@{}}
& \textbf{Set} & \textbf{open} & \textbf{closed} & \textbf{compact} \\
\hline
1 & $E_1 = \{0 < x_2 \leq 2\} \subseteq \mathbb{R}^3$ & $\times$ & $\times$ & $\times$ \\
2 & $E_2 = \{\sin|x| \geq \tfrac14\}$ & $\times$ & $\checkmark$ & $\times$ \\
3 & $E_3$ (union of open sets) & $\checkmark$ & $\times$ & $\times$ \\
4 & $E_4 = \{x\cdot y > \tfrac12|x||y|\}$ & $\checkmark$ & $\times$ & $\times$ \\
5 & $E_5 = \{X \text{ invertible}\}$ & $\checkmark$ & $\times$ & $\times$ \\
6 & $E_6 = \{X \text{ symmetric}\}$ & $\times$ & $\checkmark$ & $\times$ \\
7 & $E_7 = \{X \text{ with entries } 0,1\}$ & $\times$ & $\checkmark$ & $\checkmark$ \\
8 & $E_8 = [-6,6]^n$ & $\times$ & $\checkmark$ & $\checkmark$ \\
9 & $E_9 = E_8 \times E_3$ & $\times$ & $\times$ & $\times$ \\
10 & $E_{10} = E_2 \cap E_8$ & $\times$ & $\checkmark$ & $\checkmark$ \\
\end{tabular}
\end{center}

\begin{enumerate}[label=\textbf{(\alph*)}]
  \item Half-open in the $x_2$ direction, so \textbf{neither}: no ball around a point with
    $x_2 = 2$ stays inside, and the sequence $(0,\tfrac1k,0) \to 0 \notin E_1$ shows it is not
    closed. Unbounded, hence not compact.
  \item \textbf{Closed:} $E_2 = h^{-1}\big([\tfrac14,\infty)\big)$ for the continuous
    $h(x) := \sin(|x|)$. Not open, since a point with $\sin|x| = \tfrac14$
    has points of $E_2^c$ arbitrarily close. Not compact: $\sin$ is periodic, so $E_2$ contains
    infinitely many concentric annuli and is unbounded.
  \item \textbf{Open:} each set in the union is defined by two \emph{strict} inequalities between
    continuous functions, hence open, and an arbitrary union of open sets is open
    (\cref{item:arbitrary_union_open}). It is non-empty (take $x = (2,0,0)$ and $n=1$) and
    unbounded (take $x = (-M,0,0)$ for large $M$), so it is not compact; and a non-empty proper
    open subset of the connected space $\mathbb{R}^3$ cannot also be closed
    (\cref{def:connectedness}).
  \item \textbf{Open:} $E_4 = F^{-1}\big((0,\infty)\big)$ for the continuous
    $F(x,y) := x\cdot y - \tfrac12|x||y|$. Unbounded (scale any element),
    hence not compact, and not closed by the connectedness argument of part \textbf{(c)}.
    Geometrically $E_4$ is the set of pairs enclosing an angle strictly less than $60^\circ$
    (\cref{rem:inner_product_angle_geometry}).
  \item \textbf{Open:} $E_5 = \det^{-1}\big(\mathbb{R}\setminus\{0\}\big)$ and $\det$ is a
    polynomial in the entries, hence continuous. Unbounded, so not compact.
  \item \textbf{Closed:} $E_6$ is the solution set of the linear equations $X_{ij} - X_{ji} = 0$,
    i.e.\ the preimage of the closed set $\{0\}$ under a continuous (indeed linear) map. Being a
    linear subspace it is unbounded, hence not compact, and for $n \geq 2$ it is a proper
    subspace, hence has empty interior and is not open.
  \item \textbf{Closed and compact:} $E_7$ is \emph{finite}, with $2^{n^2}$ elements. Finite sets
    are compact (\cref{ex:compact_finite_complete}\textbf{(a)}) and closed. Not open, since no
    ball around a matrix of $0$s and $1$s consists only of such matrices.
  \item \textbf{Closed and compact:} $E_8 = [-6,6]^n$ is closed (a finite intersection of the
    closed sets $\{|x_i| \leq 6\}$) and bounded, so \cref{thm:heine_borel} applies. Not open:
    no ball around a point with some $|x_i| = 6$ fits inside.
  \item \textbf{Neither, and not compact.} For non-empty $A, B$, the product $A \times B$ is open
    exactly when both factors are, closed exactly when both are, and compact exactly when both
    are. Here $E_8$ is not open and $E_3$ is neither closed nor compact, so $E_9$ fails all
    three.

    \begin{ainote}
    The official solution opens this part with \qt{Since $E_8$ is unbounded so is $E_9$}. That
    is a slip, since $E_8 = [-6,6]^n$ is bounded. It is $E_3$ that is unbounded, as the official
    solution itself established two parts earlier. The conclusion (that $E_9$ is not compact) is
    correct; only the reason is misattributed.
    \end{ainote}
  \item \textbf{Closed and compact:} an intersection of the closed set $E_2$ with the closed
    set $E_8$ is closed (\cref{prop:open_closed_operations}\textbf{(c)}), and it is bounded
    because $E_8$ is. Heine--Borel again. Alternatively, quote
    \cref{prop:closed_subset_of_compact}: a closed subset of the compact $E_8$ is compact.
\end{enumerate}

\begin{remark}[The efficient method, extracted]
Nine of the ten sets were settled without producing a single $\varepsilon$. The recipe is:
\emph{strict} inequalities between continuous functions give open sets, \emph{weak} ones give
closed sets, and mixing the two gives neither. Compactness then reduces to checking boundedness,
by Heine--Borel. This is what the problem means by proving assertions \qt{efficiently}, and it is
the practical payoff of \cref{def:continuous}\textbf{(c)} over the $\varepsilon$--$\delta$
formulation.
\end{remark}
\end{exercisesolution}

\begin{exercisesolution}[Solution to \cref{ex:open_complete_compact_multiple_choice}]
\label{sol:open_complete_compact_multiple_choice}
% Generator: Claude Opus 5 (High)
\textbf{(a), (b), (c) and (e) are true; (d) is false.}

The engine behind \textbf{(a)}, \textbf{(b)}, \textbf{(c)} and \textbf{(e)} is a single fact,
\cref{rem:complete_vs_closed}: inside a \emph{complete} ambient space, a subset is complete if and
only if it is closed. Here the ambient space is $\mathbb{R}^n$, which is complete by
\cref{prop:Rn_coordinatewise_complete}.

\textbf{(a) True.} A compact subset of $\mathbb{R}^n$ is closed and bounded
(\cref{thm:heine_borel}). If it were also open and non-empty, it would be a non-empty
\emph{clopen} proper subset of $\mathbb{R}^n$, which is impossible because $\mathbb{R}^n$ is
connected. Indeed $\mathbb{R}^n$ is path-connected, any two points being joined by a segment, and
hence connected by \cref{lem:path_connected_implies_connected}.

\textbf{(b) True,} by the same argument: complete $\Rightarrow$ closed, and non-empty, open,
closed and proper is again impossible in the connected space $\mathbb{R}^n$.

\textbf{(c) True.} Complete $\Rightarrow$ closed, and a closed set contains all its accumulation
points. Indeed, if $x$ is an accumulation point of $A$ then some sequence in $A$ converges to $x$,
and \cref{prop:open_closed_via_sequences}\textbf{(b)} puts $x$ in $A$.

\textbf{(d) False.} Every singleton $\{q\}$ is complete (its only Cauchy sequences are
eventually constant), yet
\[\mathbb{Q} = \bigcup_{q \in \mathbb{Q}}\{q\}\]
is a countable union of complete sets which is famously \emph{not} complete
(\cref{rem:complete_vs_closed}). Equivalently, in terms of closedness: a countable union of
closed sets need not be closed, exactly as in \cref{rem:finiteness_is_necessary}.

\textbf{(e) True,} this being \cref{rem:complete_vs_closed} stated for the ambient space
$\mathbb{R}^n$, and it is worth stressing that the \qt{if} direction is where completeness of the
ambient space is used. Inside $\mathbb{Q}$ instead of $\mathbb{R}$ the statement collapses:
$\mathbb{Q}$ is closed in itself but not complete.
\end{exercisesolution}

% Reclassified from ainote to remark: this is a second proof, which is content.
\begin{remark}[A second route to \textbf{(b)}, avoiding connectedness]
Part \textbf{(b)} is also script Exercise 9.76, and the script proves it without invoking
connectedness at all. The argument is more elementary and worth recording. Fix $x \in U$ and
$y \notin U$, and let $\gamma(t) := x+t(y-x)$, $t\in[0,1]$, be the segment between them. Put
$t_* := \sup\{t\in[0,1] : \gamma(t)\in U\}$. The point $\gamma(t_*)$ cannot lie in $U$, since
openness of $U$ would let $t_*$ be pushed further out, and yet it is a limit of points
$\gamma(t_k) \in U$ with $t_k \nearrow t_*$. That is a convergent, hence Cauchy, sequence in $U$
whose limit escapes $U$.

This avoids \cref{lem:path_connected_implies_connected} entirely, at the cost of only working in
$\mathbb{R}^n$, since the segment needs the vector space structure. The connectedness proof above
generalises to any connected, locally connected ambient space.
\end{remark}

\begin{exercisesolution}[Solution to \cref{ex:uniform_lipschitz_family}]
\label{sol:uniform_lipschitz_family}
% Generator: Claude Opus 5 (High)
\begin{enumerate}[label=\textbf{(\alph*)}]
  \item By the mean value theorem, for any $x,y \in [0,1]$ there is $\xi$ between them with
    \[\lvert f_k(x)-f_k(y)\rvert = \lvert f_k'(\xi)\rvert\,\lvert x-y\rvert \leq \lvert x-y\rvert,\]
    so every $f_k$ is $1$-Lipschitz \emph{with the same constant}. Given $\varepsilon>0$, the
    choice $\delta := \varepsilon$ then satisfies \cref{def:equicontinuous} for every $k$ and
    every pair $x,y$ at once. A uniformly Lipschitz family is always equicontinuous, and this is
    by far the most common way the hypothesis is verified in practice.
  \item Apply the same bound with $y := 0$ and use $f_k(0)=0$:
    \[\lvert f_k(x)\rvert = \lvert f_k(x)-f_k(0)\rvert \leq \lvert x\rvert \leq 1
      \quad\text{for all } x \in [0,1],\]
    so $\lVert f_k\rVert_\infty \leq 1$ for every $k$. Note that the normalisation $f_k(0)=0$ is
    doing real work: without it the family $f_k :\equiv k$ satisfies the derivative bound and is
    unbounded.
  \item Both hypotheses of \cref{thm:ascoli_arzela} hold on the compact set $[0,1]$, so some
    subsequence converges uniformly to a continuous limit.

    Replacing the uniform bound by mere differentiability destroys this, and
    \cref{ex:bounded_no_convergent_subsequence} is already the counterexample: $f_k(x) := x^k$
    is $C^1$ on $[0,1]$ with $f_k(0)=0$, so it satisfies every hypothesis \emph{except} the
    uniform one. Its derivative $f_k'(x) = kx^{k-1}$ has $\lVert f_k'\rVert_\infty = k$, which
    grows without bound. And we saw that no subsequence of $(x^k)$ converges. What matters is
    therefore not that each $f_k$ is well behaved, but that they are well behaved \emph{by a
    common margin}.
\end{enumerate}
\end{exercisesolution}

\begin{exercisesolution}[Proof of \cref{ex:annulus_continuous_image}]
\label{sol:annulus_continuous_image}
% Generator: Gemini 3.6 Flash (High)
% Correction: Claude Opus 5 (High) --- solution for part (c) supplied; U is bounded by
% \overline{B_2(0)}, not B_2(0), since (0,2) lies in U.
\begin{enumerate}[label=\textbf{(\alph*)}]
  \item The map $(x,y) \mapsto x^2+y^2$ is continuous, and $U$ is the preimage of the closed set
    $[1,4]$ under it, so $U$ is closed. It is also bounded, being contained in
    $\overline{B_2(0)}$. By Heine--Borel (\cref{thm:heine_borel}), $U$ is compact. It is moreover
    path-connected --- any two of its points can be joined by moving along a circle and then
    radially --- and therefore connected.

  \item Since $f$ is continuous and $U$ is compact, $f(U)$ is compact
    (\cref{item:continuity_preserves_compact}); since $U$ is connected, $f(U)$ is connected
    (\cref{thm:continuity_preserves_connected}). The connected subsets of $\mathbb{R}$ are exactly
    the intervals, and a compact interval is one of the form $[a,b]$. It remains to rule out
    $a = b$, that is, to check $f$ is non-constant: $(1,0)$ and $(0,2)$ both lie in $U$, and
    $f(1,0) = \sin 0 - 0 = 0$ while $f(0,2) = \sin 0 - 16 = -16$. Hence $f(U) = [a,b]$ with $a < b$.

  \item \textbf{No.} A loop running once around the inner hole cannot be contracted to a point
    within $U$ without crossing the removed disc $\{x^2+y^2 < 1\}$. The annulus is the standard
    example separating \qt{connected} from \qt{simply connected}, and it is worth keeping the two
    apart: part \textbf{(a)} is what \textbf{(b)} needs, and simple connectedness plays no role
    there at all.
\end{enumerate}
\end{exercisesolution}

\begin{exercisesolution}[Proof of \cref{ex:nearest_point_to_origin}]
\label{sol:nearest_point_to_origin}
% Generator: Claude Opus 5 (High)
\begin{enumerate}[label=\textbf{(\alph*)}]
  \item The set $A$ need not be bounded, so Weierstrass does not apply to it directly. The remedy is
    to cut it down to a compact piece without losing the minimum. Fix any $a_0 \in A$ and set
    $R := \lvert a_0\rvert$. The set
\[K := A \cap \overline{B_R(0)}\]
    is closed (an intersection of two closed sets) and bounded, hence compact by Heine--Borel
    (\cref{thm:heine_borel}), and it is non-empty because $a_0 \in K$. Since $f$ is continuous, it
    attains a minimum on $K$ (\cref{item:extreme_value_theorem}), say at $P_0 \in K \subseteq A$.
    This $P_0$ minimises over all of $A$: any $P \in A \setminus K$ has $\lvert P\rvert > R =
    \lvert a_0\rvert \ge f(P_0)$, because $a_0 \in K$.

  \item Take $A := \{x \in \mathbb{R}^n \mid 0 < \lvert x\rvert \le 1\}$, the closed unit ball with
    its centre removed. Then $\inf_A f = 0$, since $\lvert x\rvert$ can be made as small as we like,
    but $f(x) = 0$ would force $x = 0 \notin A$. So the infimum is not attained. The set is bounded,
    which shows that closedness --- not boundedness --- is the hypothesis doing the work in
    \textbf{(a)}.

  \item Let $A$ be convex and suppose $P_0, P_1 \in A$ both minimise, with common value
    $m := \lvert P_0\rvert = \lvert P_1\rvert$. If $m = 0$ then $P_0 = P_1 = 0$ and there is nothing
    to prove, so assume $m > 0$. By convexity the midpoint $M := \tfrac12(P_0+P_1)$ lies in $A$, and
    the triangle inequality gives
\[\lvert M\rvert \le \tfrac12\big(\lvert P_0\rvert + \lvert P_1\rvert\big) = m.\]
    Equality in the triangle inequality for the Euclidean norm forces $P_0$ and $P_1$ to be
    non-negative multiples of one another; having equal non-zero norms, they would then be equal.
    So if $P_0 \neq P_1$ the inequality is strict, $\lvert M\rvert < m$, contradicting minimality of
    $m$. Hence $P_0 = P_1$.

  \item Let $P_0 \in A$ be a minimiser and put $m := \lvert P_0\rvert$, which is positive because
    $0 \notin A$. Suppose, for contradiction, that $P_0$ were an interior point of $A$, so that
    $B_\varepsilon(P_0) \subseteq A$ for some $\varepsilon > 0$. Note first that
    $\varepsilon \le m$: otherwise $\lvert 0 - P_0\rvert = m < \varepsilon$ would put the origin
    inside $B_\varepsilon(P_0) \subseteq A$, contrary to $0 \notin A$. Now move from $P_0$ straight
    towards the origin: the point
\[P_1 := \left(1 - \frac{\varepsilon}{2m}\right) P_0\]
    satisfies $\lvert P_1 - P_0\rvert = \frac{\varepsilon}{2m}\lvert P_0\rvert = \frac{\varepsilon}{2}
    < \varepsilon$, so $P_1 \in A$. Since $\varepsilon \le m$, the scalar $1 - \frac{\varepsilon}{2m}$
    lies in $[\tfrac12, 1)$ and in particular is positive, so
    $\lvert P_1\rvert = \left(1 - \frac{\varepsilon}{2m}\right)m = m - \frac{\varepsilon}{2} < m$.
    This contradicts the minimality of $\lvert P_0\rvert$. Therefore $P_0$ is not interior, and since
    $P_0 \in A \subseteq \overline{A}$ it must lie in $\partial A$.
% Correction: Claude Opus 5 (High) --- the bound epsilon <= m was implicit in the first draft. It is
% needed: without it the scalar could be negative and |P_1| would be epsilon/2 - m, not m -
% epsilon/2. The official solution secures the same point by taking epsilon := min{epsilon', |P_0|}.
\end{enumerate}
\begin{remark}
The four parts isolate four different hypotheses and show what each one buys: closedness gives
\emph{existence}, convexity gives \emph{uniqueness}, and $0 \notin A$ gives \emph{location}. Part
\textbf{(b)} is what stops the list reading as decoration. It is worth noticing that boundedness is
never assumed: the trick in \textbf{(a)} of intersecting with a ball large enough to catch one
point of $A$ is the standard way to run a Weierstrass argument on an unbounded closed set, and it
recurs whenever a coercive function is minimised.
\end{remark}
\end{exercisesolution}

% Generator: Claude Opus 5 (High)
% No solution key ships with the Jossen papers, so this is an independent derivation.
\begin{exercisesolution}[\cref{ex:pseudocompactness_implications}]
\label{sol:pseudocompactness_implications}
\begin{enumerate}[label=\textbf{(\alph*)}]
  \item Assume \textbf{(ii)} and let $f : X \to \mathbb{R}$ be continuous. It attains a maximum
    at some $p$ and a minimum at some $q$, so $f(X) \subseteq [f(q),f(p)]$ and $f$ is bounded.

    Conversely, assume \textbf{(i)} and let $f : X \to \mathbb{R}$ be continuous. Then
    $M := \sup f(X)$ is a real number, since $f(X)$ is non-empty and bounded. Suppose $f$ never
    took the value $M$. Then $M - f(x) > 0$ for every $x$, so
    \[g : X \to \mathbb{R}, \qquad g(x) := \frac{1}{M-f(x)}\]
    is well defined and continuous. By definition of the supremum there is a sequence $(x_k)$ in
    $X$ with $f(x_k) \to M$, whence $g(x_k) \to +\infty$. Therefore $g$ is an unbounded continuous
    function, contradicting \textbf{(i)}. Consequently $f$ attains its maximum, and applying this
    to $-f$ gives the minimum.

  \item Assume \textbf{(iii)} and let $f : X \to \mathbb{R}$ be continuous. The preimages
    $f^{-1}\big((-k,k)\big)$, $k \in \mathbb{N}$, are open (\cref{def:continuous}) and cover $X$,
    because every $f(x)$ is a real number and so lies in $(-k,k)$ for $k$ large. Compactness
    yields finitely many indices $k_1 < \dots < k_r$ whose preimages already cover $X$, and these
    are nested, so $X = f^{-1}\big((-k_r,k_r)\big)$. Hence $\lvert f\rvert < k_r$ throughout.

  \item Assume \textbf{(iv)} and suppose some continuous $f : X \to \mathbb{R}$ were unbounded.
    Choose $x_k \in X$ with $\lvert f(x_k)\rvert \ge k$ for every $k$. By \textbf{(iv)} the
    sequence $(x_k)$ has an accumulation point $x \in X$, so some subsequence satisfies
    $x_{k_j} \to x$. Continuity gives $f(x_{k_j}) \to f(x)$, so the numbers $f(x_{k_j})$ are
    bounded. But $\lvert f(x_{k_j})\rvert \ge k_j \to \infty$, a contradiction.
\end{enumerate}
\begin{remark}[Weierstrass, split into its two halves]
Part \textbf{(a)} is the striking one, and it is worth seeing what it says: the two conclusions of
the extreme value theorem, \emph{bounded} and \emph{attained}, are not independent facts about the
space. Either one forces the other, with no compactness assumed anywhere. What compactness (or
\textbf{(iv)}) supplies is only the entry ticket \textbf{(i)}.

The trick in \textbf{(a)} is worth keeping: to force a supremum to be attained, divide by the gap
$M - f$ and let the failure to attain it manufacture an unbounded function. It reappears whenever a
statement about suprema has to be converted into a statement about boundedness.

None of the reverse implications hold. A space satisfying \textbf{(i)} is called
\newterm{pseudocompact}, and there are pseudocompact metric spaces that are not compact. Inside
$\mathbb{R}^n$ the gap closes, by \cref{thm:heine_borel}. The chain proved here is
\textbf{(iii)} $\Rightarrow$ \textbf{(i)} $\Leftrightarrow$ \textbf{(ii)} and \textbf{(iv)}
$\Rightarrow$ \textbf{(i)}; that \textbf{(iii)} and \textbf{(iv)} are themselves equivalent for
metric spaces is \cref{thm:sequential_topological_compactness_equivalence}, which is a much deeper
statement than either implication above.
\end{remark}
\end{exercisesolution}

% Generator: Claude Opus 5 (High)
\begin{exercisesolution}[\cref{ex:sumset_open_closed_compact}]
\label{sol:sumset_open_closed_compact}
\begin{enumerate}[label=\textbf{(\alph*)}]
  \item Suppose $A$ is open. For fixed $b \in \mathbb{R}^2$ the translate
    $A + b := \{a+b \mid a \in A\}$ is open, since translation is a homeomorphism of
    $\mathbb{R}^2$; concretely, $B_r(a) \subseteq A$ gives $B_r(a+b) \subseteq A+b$. Now
    \[A + B = \bigcup_{b \in B} (A+b)\]
    is a union of open sets and therefore open. The case of $B$ open is symmetric, since
    $A+B = B+A$.

  \item Let $c \in \overline{A+B}$ and pick $c_k = a_k + b_k \in A+B$ with $c_k \to c$, where
    $a_k \in A$ and $b_k \in B$. Since $B$ is compact, it is sequentially compact
    (\cref{thm:sequential_topological_compactness_equivalence}), so a subsequence
    $b_{k_j} \to b \in B$ converges. Then
    \[a_{k_j} = c_{k_j} - b_{k_j} \longrightarrow c - b,\]
    and $A$ is closed, so $c - b =: a$ lies in $A$. Hence $c = a + b \in A+B$, and $A+B$ is closed.

  \item The product $A \times B \subseteq \mathbb{R}^4$ is compact: it is closed and bounded,
    because $A$ and $B$ are, so \cref{thm:heine_borel} applies. Addition
    \[s : \mathbb{R}^2\times\mathbb{R}^2 \to \mathbb{R}^2, \qquad s(a,b) := a+b\]
    is continuous, being linear, and $A+B = s(A\times B)$ is the continuous image of a compact set,
    hence compact.
\end{enumerate}
\begin{remark}[Compactness of $B$ is not decoration]
Two closed sets can have a non-closed sum, so \textbf{(b)} genuinely needs more than closedness.
Take
\[A := \{(n,0) \mid n \in \mathbb{N}\}, \qquad B := \{(-n,\tfrac1n) \mid n \in \mathbb{N}\}.\]
Both are closed, since each is a set of isolated points running off to infinity and so has no
accumulation point at all. Their sum contains $(0,\tfrac1n)$ for every $n$, and those converge to
the origin. But $(0,0) \in A+B$ would need $n = m$ and $\tfrac1m = 0$ at once, so the origin is
missing and $A+B$ is not closed.

What compactness of $B$ supplies in the proof is precisely the convergent subsequence of
$(b_k)$: closedness of $B$ alone would let $(b_k)$ escape to infinity, which is exactly what the
counterexample arranges. Note also that nothing in any of the three arguments used the dimension,
so all three hold verbatim in $\mathbb{R}^n$, and \textbf{(a)} and \textbf{(b)} hold in any normed
space.
\end{remark}
\end{exercisesolution}

\begin{exercisesolution}[Solution to \cref{ex:kobel_compactness_characterisation}]
\label{sol:kobel_compactness_characterisation}
% Generator: Claude Opus 5 (High)
Only \textbf{(c)} is true.
\begin{enumerate}[label=\textbf{(\alph*)}]
  \item \textbf{False.} This is Heine--Borel (\cref{thm:heine_borel}), and it is a theorem about
    $\mathbb{R}^n$ with the standard metric, not about metric spaces. \Cref{ex:heine_borel_fails}
    breaks it with the bounded metric $\tilde d$, under which $\mathbb{N}$ is closed and bounded
    and not compact; an infinite set with the discrete metric does the same.
  \item \textbf{False.} Being open and closed is a statement about how $E$ sits in $X$, and it has
    no bearing on compactness. Any space is open and closed in itself, so $E = X = \mathbb{R}$ is
    a counterexample.
  \item \textbf{True}, and it is the honest general replacement for \textbf{(a)}: a metric space is
    compact if and only if it is complete and totally bounded
    (\cref{rem:compact_iff_complete_totally_bounded}, built on
    \cref{lem:sequentially_compact_totally_bounded}).
  \item \textbf{False.} Total boundedness alone is not enough --- completeness is the other half of
    \textbf{(c)} --- and openness does not supply it. The interval $(0,1) \subseteq \mathbb{R}$ is
    open and totally bounded and not compact.
\end{enumerate}
\begin{remark}[Why \textbf{(a)} is the tempting one]
\Cref{rem:boundedness_not_topological} already contains the reason \textbf{(a)} fails and
\textbf{(c)} succeeds, and it is worth restating in one line, because the question is designed to
catch a reader who remembers Heine--Borel without its hypothesis. Compactness is a topological
property; boundedness is not, since two metrics with the same open sets can disagree about which
sets are bounded. Total boundedness, unlike boundedness, is not so easily destroyed by rescaling
the metric, and that is exactly why it is the notion that survives into the general statement.
\end{remark}
\end{exercisesolution}

\begin{exercisesolution}[Solution to \cref{ex:box_preimage_of_compact_not_compact}]
\label{sol:box_preimage_of_compact_not_compact}
% Generator: Claude Opus 5 (High)
Take the projection onto the first coordinate,
\[g(x,y) := (x,\,0), \qquad K := \{(0,0)\}.\]
Then $g$ is linear and therefore continuous, and it is nonconstant since $g(1,0) = (1,0)$ while
$g(0,0) = (0,0)$. The set $K$ is a single point, hence compact. But
\[g^{-1}(K) = \{(x,y) \mid x = 0\} = \{0\}\times\mathbb{R}\]
is the whole $y$-axis, which is closed but unbounded and so not compact by
\cref{thm:heine_borel}.

\begin{remark}[Continuity preserves compactness in one direction only]
\Cref{item:continuity_preserves_compact} says that $g(C)$ is compact whenever $C$ is, and it is easy
to carry that away as a statement about $g$ and compactness rather than about images. It is not:
the preimage of a compact set under a continuous map can be as large as one likes, and the
example above makes it a whole line by collapsing one direction. A map that \emph{does} have the
property is called \newterm{proper}, and properness is a genuine extra hypothesis, not a
consequence of continuity. \Cref{ex:box_proper_map_not_diffeomorphism} constructs one, and shows
that properness in turn buys much less than one might hope.

The official key gives a different answer with the same moral, $g(x,y) := (\arctan x, \arctan y)$
together with any compact $K$ containing the open square $(-\tfrac\pi2,\tfrac\pi2)^2$, for which
$g^{-1}(K) = \mathbb{R}^2$. There the map is injective; it is the \emph{boundedness} of its
image that does the work, rather than a collapsed direction.
\end{remark}
\end{exercisesolution}


